% formal/product-two-bounce-class.tex
%
% Developer-facing proof packet for the two-bounce class of a planar
% Lagrangian product.  All mathematics in the red status frame is
% agent-derived and awaits Jörn's review.

\subsection{Status, scope, and source conventions}
\label{subsec:product-two-bounce-status}

This packet separates three claims which must not be conflated:
\begin{enumerate}
  \item the geometric minimum among strong two-bounce Minkowski billiards;
  \item the optimum of the complete mathematical \(k=2\) product QP family,
    before transition pruning or numerical solution;
  \item the value returned from the transition-pruned, f64-seeded candidate
    stream used by
    \path{experiments/sys-datascience/methods/product-bounce-distribution/src/bin/class-minima.rs}.
\end{enumerate}
The first two values are proved below to equal the same difference-body
inradius.  Equality for the implementation stream is conditional on retaining
and exactly resolving one equality candidate.  Current sources do not prove
that coverage condition universally.

The strong-billiard convention and its two normal-cone inclusions are
Rudolf's published Definition~2 \cite[Definition~2]{Rudolf2022}.  The
published Theorem~1 identifies the unrestricted strong-billiard minimum with
the EHZ capacity, but that global theorem is not used to prove the
class-restricted statements below.  The QP orientation and the action formula
\(1/(2Q_\sigma)\) are those of
\texttt{thesis/04-haim-kislev-quadratic-program.tex}.

\begin{unverified}
All new mathematical claims in this packet have been checked by agents and
compiled, but have not been accepted by Jörn.  In particular, \emph{proved}
below means proved within this packet at agent-review status, not approved for
publication in the thesis.
\end{unverified}

\subsection{Conventions}
\label{subsec:product-two-bounce-conventions}

Let \(P,Q\subset\mathbb R^2\) be compact full-dimensional convex polygons.
For a convex body \(C\), write
\[
 h_C(u)=\max_{x\in C}\langle u,x\rangle,
 \qquad
 C^\circ=\{u:h_C(u)\leq1\},
\]
and use the outer normal cone
\[
 N_C(x)=\{u:\langle u,z-x\rangle\leq0
                 \text{ for every }z\in C\}.
\]
The exposed face with outer normal \(u\neq0\) is
\[
 F_C(u)=\{x\in C:\langle u,x\rangle=h_C(u)\}.
\]
Put
\[
 D_P=P-P,\qquad D_Q=Q-Q.
\]
These difference bodies are centrally symmetric convex bodies with zero in
their interiors.  The polar convention gives
\[
 \mu_{D_Q^\circ}(d)=h_{D_Q}(d).
\]

A \emph{nondegenerate strong two-bounce \((P,Q)\)-billiard} consists of
distinct \(q_1,q_2\in\partial P\) and \(p_1,p_2\in\partial Q\) such that,
with \(d=q_2-q_1\neq0\),
\begin{equation}
\label{eq:product-two-bounce-strong-relations}
 d\in N_Q(p_1),\qquad
 -d\in N_Q(p_2),\qquad
 p_2-p_1\in-N_P(q_2),\qquad
 p_1-p_2\in-N_P(q_1).
\end{equation}
Its closed length in the repository convention is
\begin{equation}
\label{eq:product-two-bounce-length}
 \ell_Q(q_1,q_2)
 =h_Q(d)+h_Q(-d)
 =h_{D_Q}(d).
\end{equation}
There is no additional factor of two.  If \(Q=-Q\), then \(D_Q=2Q\), and
only in that special case does the last expression become \(2h_Q(d)\).

For the QP discussion, translate \(P\) and \(Q\), without changing their
difference bodies or the billiard length, so that
\(0\in\operatorname{int}P\cap\operatorname{int}Q\).  Write irredundant facet
representations
\[
 P=\{q:\langle a_i,q\rangle\leq1\},
 \qquad
 Q=\{p:\langle b_j,p\rangle\leq1\}.
\]
The \(a_i\) are the vertices of \(P^\circ\), and the \(b_j\) are the
vertices of \(Q^\circ\).  Product coordinates have the project order
\((q_1,q_2,p_1,p_2)\), and
\[
 \omega_0((u,0),(0,v))=\langle u,v\rangle.
\]

\subsection{Two-point translatability and the difference-body value}
\label{subsec:product-two-bounce-difference-body}

\begin{unverified}
\begin{lemma}[Two-point translatability]
\label{lem:product-two-bounce-translatability}
For \(d\in\mathbb R^2\), a two-point set with displacement \(d\) is
translatable into \(\operatorname{int}P\) if and only if
\[
 d\in\operatorname{int}(P-P).
\]
\end{lemma}

\begin{proof}
Translation does not change displacement.  Thus the asserted equivalence is
the identity
\[
 \operatorname{int}P-\operatorname{int}P
 =\operatorname{int}(P-P).
\]
The left side is open and contained in \(P-P\), hence is contained in the
right side.  Conversely, take \(d\in\operatorname{int}(P-P)\).  If \(d=0\),
write \(d=p_0-p_0\) with \(p_0\in\operatorname{int}P\).  If \(d\neq0\),
the ray through \(d\) continues a little past \(d\) inside \(P-P\); hence
\(\lambda d=x-y\) for some \(\lambda>1\) and \(x,y\in P\).  For
\(p_0\in\operatorname{int}P\), the points
\[
 x'=\lambda^{-1}x+(1-\lambda^{-1})p_0,\qquad
 y'=\lambda^{-1}y+(1-\lambda^{-1})p_0
\]
belong to \(\operatorname{int}P\) and satisfy \(x'-y'=d\).
\end{proof}

\begin{definition}[Difference-body two-bounce value]
\label{def:product-two-bounce-a2}
Define
\begin{equation}
\label{eq:product-two-bounce-boundary}
 A_2(P,Q)=\min_{d\in\partial D_P}h_{D_Q}(d).
\end{equation}
\end{definition}

\begin{lemma}[Boundary and containment formulations]
\label{lem:product-two-bounce-containment}
The minimum \(Q\)-length of a non-translatable two-point polygon is
\(A_2(P,Q)\), and
\begin{equation}
\label{eq:product-two-bounce-containment}
 A_2(P,Q)
 =\max\{r\geq0:rD_Q^\circ\subseteq D_P\}.
\end{equation}
The minimum and maximum are positive and attained.
\end{lemma}

\begin{proof}
By Lemma~\ref{lem:product-two-bounce-translatability}, the admissible
displacements are \(\mathbb R^2\setminus\operatorname{int}D_P\).
The function \(h_{D_Q}\) is positive away from zero and positively
homogeneous.  On each ray, the first admissible point is the unique point of
\(\partial D_P\) on that ray; every point farther out has at least its
\(h_{D_Q}\)-value.  Compactness of \(\partial D_P\) gives
\eqref{eq:product-two-bounce-boundary} and attainment.

Set \(C=D_Q^\circ\), so that \(h_{D_Q}=\mu_C\).  For two convex bodies
\(B,C\) containing zero in their interiors, radial comparison gives
\[
 \max\{r:rC\subseteq B\}=\min_{x\in\partial B}\mu_C(x).
\]
Indeed, \(rC\subseteq B\) precisely when the radial size of \(rC\) does not
exceed that of \(B\) on any ray, which is equivalent to
\(\mu_C(x)\geq r\) for every \(x\in\partial B\).
Taking \(B=D_P\) proves \eqref{eq:product-two-bounce-containment}.
Positivity follows because both bodies contain zero in their interiors.
\end{proof}
\end{unverified}

\subsection{Strong two-bounce realization}
\label{subsec:product-two-bounce-strong-realization}

The next theorem proves the optimizer correspondence directly.  It is
stronger than equality of the two minimum values: every minimizing
displacement is realized by a strong billiard.

\begin{unverified}
\begin{theorem}[Strong two-bounce formula]
\label{thm:product-strong-two-bounce}
For full-dimensional convex polygons \(P,Q\subset\mathbb R^2\),
\begin{equation}
\label{eq:product-strong-two-bounce-minimum}
 \min\{\ell_Q(q_1,q_2):
       (q_1,q_2;p_1,p_2)\text{ is a nondegenerate strong two-bounce
       billiard}\}
 =A_2(P,Q).
\end{equation}
Moreover, every \(d\in\partial D_P\) satisfying
\(h_{D_Q}(d)=A_2(P,Q)\) occurs as \(q_2-q_1\) in a minimizing strong
two-bounce billiard.
\end{theorem}

\begin{proof}
Put \(r=A_2(P,Q)\), \(C=D_Q^\circ\), and choose a minimizing
\(d\in\partial D_P\).  Then \(h_{D_Q}(d)=r\), so
\(x=d/r\in\partial C\).  Lemma~\ref{lem:product-two-bounce-containment}
gives
\[
 rC\subseteq D_P,\qquad d\in\partial(rC)\cap\partial D_P.
\]
Choose \(0\neq u\in N_{D_P}(d)\).  Because \(d\in rC\subseteq D_P\),
\[
 h_{rC}(u)\geq\langle u,d\rangle
 =h_{D_P}(u)\geq h_{rC}(u).
\]
Thus equality holds throughout, and \(u\in N_{rC}(d)=N_C(x)\).

For Minkowski differences,
\begin{equation}
\label{eq:product-two-bounce-face-decompositions}
 F_{D_P}(u)=F_P(u)-F_P(-u),\qquad
 F_{D_Q}(x)=F_Q(x)-F_Q(-x).
\end{equation}
Since \(d\in F_{D_P}(u)\), choose
\[
 q_2\in F_P(u),\qquad q_1\in F_P(-u),\qquad d=q_2-q_1.
\]
Both points lie on \(\partial P\), and they are distinct because \(d\neq0\).

The polar body is the sublevel set
\(C=\{z:h_{D_Q}(z)\leq1\}\).  At \(x\in\partial C\), the normal-cone
formula for a sublevel set of a support function is
\[
 N_C(x)=\{\lambda y:\lambda\geq0,\ y\in F_{D_Q}(x)\}.
\]
Consequently, \(u=\lambda y\) for some \(\lambda>0\) and
\(y\in F_{D_Q}(x)\).  By
\eqref{eq:product-two-bounce-face-decompositions}, choose
\[
 p_1\in F_Q(x),\qquad p_2\in F_Q(-x),\qquad y=p_1-p_2.
\]
Because \(d=rx\), these exposed-face conditions give
\[
 d\in N_Q(p_1),\qquad -d\in N_Q(p_2).
\]
Because \(u\) is a positive multiple of \(y\), the choices of \(q_1,q_2\)
give
\[
 y\in N_P(q_2),\qquad -y\in N_P(q_1).
\]
Equivalently,
\[
 p_2-p_1=-y\in-N_P(q_2),\qquad
 p_1-p_2=y\in-N_P(q_1).
\]
These are all four relations in
\eqref{eq:product-two-bounce-strong-relations}.  The resulting billiard has
length \(h_{D_Q}(d)=r\).

Conversely, let a nondegenerate strong two-bounce billiard be given, and set
\[
 d=q_2-q_1,\qquad y=p_1-p_2.
\]
The vector \(y\) is nonzero.  Otherwise, the first two strong relations put
both \(d\) and \(-d\) in the normal cone of \(Q\) at one point.  A normal
cone of a full-dimensional convex body is pointed, so this would force
\(d=0\).
The last two strong relations give
\[
 y\in N_P(q_2),\qquad -y\in N_P(q_1).
\]
Therefore
\[
 \langle y,d\rangle
 =h_P(y)+h_P(-y)
 =h_{D_P}(y)>0,
\]
and \(d=q_2-q_1\in F_{D_P}(y)\subset\partial D_P\).  Hence
\[
 \ell_Q(q_1,q_2)=h_{D_Q}(d)\geq A_2(P,Q).
\]
Together with the construction, this proves
\eqref{eq:product-strong-two-bounce-minimum} and realization of every
minimizing displacement.
\end{proof}
\end{unverified}

Thus the minimizing displacement classes and the minimizing strong
two-bounce classes correspond: every strong two-bounce billiard supplies a
boundary displacement, and every minimizing boundary displacement has at
least one strong realization.  This is not a setwise identification of all
non-translatable pairs with all billiards.  A general
\(d\notin\operatorname{int}D_P\) need not satisfy the normal-cone relations,
and even a minimizing \(d\) can have several realizations through different
face decompositions.

\subsection{The complete mathematical two-bounce product QP}
\label{subsec:product-two-bounce-complete-qp}

A broad two-block word has cyclic form
\[
 \mathcal Q_+\,\mathcal P_+\,\mathcal Q_-\,\mathcal P_-,
\]
where each block is a singleton or an ordered adjacent pair of facets of the
indicated factor.  Initially, allow the same facet to occur in the two blocks
of one factor.  Give every active occurrence a positive weight \(\beta_k\)
satisfying
\begin{equation}
\label{eq:product-two-bounce-qp-constraints}
 \sum_k\beta_k=1,\qquad
 \sum_k\beta_k\widehat a_{\sigma_k}=0,
\end{equation}
where a product \(q\)-facet has dual vertex
\(\widehat a_i=(a_i,0)\), while a product \(p\)-facet has dual vertex
\(\widehat b_j=(0,b_j)\).  The active-traversal objective is
\[
 Q_\sigma(\beta)
 =\sum_{j<k}\beta_j\beta_k\,
   \omega_0(\widehat a_{\sigma_j},\widehat a_{\sigma_k}),
\]
with \(\widehat a_{\sigma_k}\) understood to mean the appropriate
\(\widehat a_i\) or \(\widehat b_j\).

The complete mathematical repository family additionally requires distinct
facets, equivalently disjoint blocks of the same factor.  Zero weights are
deleted from an active word.  The repeated-facet repair below proves that the
distinctness requirement does not lower the \(k=2\) optimum.

\begin{unverified}
\begin{lemma}[Collapse of an alternating four-block word]
\label{lem:product-two-bounce-block-collapse}
Let \((\sigma,\beta)\) be a positive feasible broad word.  If \(u\) is the
weighted planar sum in \(\mathcal Q_+\) and \(v\) is the weighted planar sum
in \(\mathcal P_+\), factorwise closure makes the four block sums, in cyclic
order,
\[
 u,\quad v,\quad -u,\quad -v.
\]
They satisfy \(u,v\neq0\),
\begin{equation}
\label{eq:product-two-bounce-block-normalization}
 h_{D_P}(u)+h_{D_Q}(v)=1,
\end{equation}
and
\begin{equation}
\label{eq:product-two-bounce-block-objective}
 Q_\sigma(\beta)=2\langle u,v\rangle.
\end{equation}
\end{lemma}

\begin{proof}
Consider a \(q\)-facet block of total mass \(s>0\) and weighted planar sum
\(w\).  The point \(w/s\) is either a vertex of \(P^\circ\) or an interior
point of an edge joining two adjacent vertices.  Hence
\(w/s\in\partial P^\circ\), and
\[
 h_P(w)=\mu_{P^\circ}(w)=s.
\]
Applying this to the two \(q\)-facet blocks with sums \(u,-u\), and then to
the two \(p\)-facet blocks with sums \(v,-v\), turns
\(\sum_k\beta_k=1\) into
\eqref{eq:product-two-bounce-block-normalization}.  A positive combination
of one vertex or of two adjacent vertices of a polar polygon cannot vanish,
because the corresponding boundary vertex or edge does not contain zero.
Thus \(u,v\neq0\).

Pairings within one factor vanish.  The four nonzero cross-block
contributions are
\[
 \langle u,v\rangle,\quad
 \langle u,-v\rangle,\quad
 -\langle -u,v\rangle,\quad
 \langle -u,-v\rangle.
\]
Their sum is \(2\langle u,v\rangle\), proving
\eqref{eq:product-two-bounce-block-objective}.  This calculation fixes the
repository orientation.  Reversing a fixed cyclic word negates its
objective; the complete ordered family contains the reverse orientation, so
we choose the orientation with positive objective when necessary.
\end{proof}

\begin{lemma}[Radial block encoding]
\label{lem:product-two-bounce-radial-encoding}
For every \(0\neq w\in\mathbb R^2\), the vector \(w\) is a positive weighted
sum of either one vertex or two adjacent vertices of \(P^\circ\), with total
coefficient \(h_P(w)\).  The analogous statement holds for \(Q^\circ\).
\end{lemma}

\begin{proof}
The radial point \(w/h_P(w)\) belongs to \(\partial P^\circ\).  In a polygon
it is either a vertex or lies in the relative interior of an edge.  In the
first case use the single vertex.  In the second, write the radial point as a
strict convex combination of the two endpoints of that edge and multiply
the coefficients by \(h_P(w)\).
\end{proof}

\begin{lemma}[Polar pairing value]
\label{lem:product-two-bounce-polar-pairing}
Let
\[
 M(P,Q)=\max\{\langle x,y\rangle:
                    x\in D_P^\circ,\ y\in D_Q^\circ\}.
\]
Then
\begin{equation}
\label{eq:product-two-bounce-polar-pairing}
 M(P,Q)=\frac1{A_2(P,Q)}.
\end{equation}
\end{lemma}

\begin{proof}
For fixed \(y\),
\[
 \max_{x\in D_P^\circ}\langle x,y\rangle
 =h_{D_P^\circ}(y)=\mu_{D_P}(y).
\]
Therefore
\[
 M(P,Q)=\max_{y\in D_Q^\circ}\mu_{D_P}(y).
\]
For \(r\geq0\),
\[
 rD_Q^\circ\subseteq D_P
 \quad\Longleftrightarrow\quad
 r\,\mu_{D_P}(y)\leq1
 \text{ for every }y\in D_Q^\circ
 \quad\Longleftrightarrow\quad
 rM(P,Q)\leq1.
\]
The largest such \(r\) is \(A_2(P,Q)\) by
Lemma~\ref{lem:product-two-bounce-containment}, proving the identity.  This
argument also fixes which polar belongs to which factor.
\end{proof}

\begin{proposition}[Removing repeated facets at the class optimum]
\label{prop:product-two-bounce-distinct-repair}
Every maximizing broad four-block candidate can be changed, without changing
its objective, to a maximizing candidate whose two blocks in each factor are
disjoint.  After zero coefficients are deleted, all four blocks remain
nonempty and have length one or two.
\end{proposition}

\begin{proof}
Consider the radial boundary points of \(P^\circ\) in the directions \(u\)
and \(-u\) supplied by
Lemma~\ref{lem:product-two-bounce-block-collapse}.  Their minimal faces are
vertices or edges and can share at most one polar vertex.  Indeed, two
distinct polygon edges share at most one vertex, while if the two minimal
faces shared an edge, that edge would contain the two opposite radial points
and hence zero, contradicting \(0\in\operatorname{int}P^\circ\).

Suppose the only possible shared vertex \(a\) occurs with coefficients
\(c_+>0\) and \(c_->0\) in the two temporal \(q\)-blocks.  Hold
\(c=c_++c_-\) fixed and replace these coefficients by
\[
 \theta,\qquad c-\theta,\qquad 0\leq\theta\leq c,
\]
leaving every other coefficient unchanged.  Closure and normalization depend
only on the total coefficient of \(a\), so they remain fixed.  If
\(u(\theta)\) is the first \(q\)-block sum, the other is
\(-u(\theta)\), and
\[
 Q_\sigma(\theta)=2\langle u(\theta),v\rangle
\]
is affine in \(\theta\).  One endpoint is no worse than the original split
and deletes one occurrence of \(a\).  Since the original candidate is a
broad-family maximum, that endpoint has the same value.

Neither endpoint can erase an entire \(q\)-block.  If one of the original
blocks consisted only of \(a\), factorwise closure would express zero as a
positive combination of \(a\) and the one or two vertices in the opposite
block.  Because the opposite minimal face shares \(a\), all these vertices
lie in one boundary edge (or at one boundary vertex) of \(P^\circ\), which
cannot contain zero.  Thus each block has a non-shared active vertex and
survives deletion of the repeated occurrence.

This removes the only possible repeated \(q\)-facet.  Apply the identical
argument to the two \(p\)-blocks.  The result has disjoint same-factor blocks,
four nonempty blocks, and unchanged maximal objective.
\end{proof}

\begin{theorem}[Complete mathematical \(k=2\) product QP]
\label{thm:product-two-bounce-complete-qp}
Let \(Q_{2,\max}(P,Q)\) be the maximum of \(Q_\sigma\) over the complete
mathematical \(k=2\) product family: alternating cyclic words with two
nonempty \(q\)-blocks and two nonempty \(p\)-blocks, every block a singleton
or an ordered adjacent pair, same-factor blocks disjoint, and weights
satisfying \eqref{eq:product-two-bounce-qp-constraints}.  This is the family
before transition pruning or numerical proposal.  Then
\begin{equation}
\label{eq:product-two-bounce-complete-qp}
 \begin{split}
 Q_{2,\max}(P,Q)
 &=\frac12\max_{\substack{x\in D_P^\circ\\y\in D_Q^\circ}}
      \langle x,y\rangle\\
 &=\frac1{2A_2(P,Q)},
 \qquad
 \frac1{2Q_{2,\max}(P,Q)}=A_2(P,Q).
 \end{split}
\end{equation}
In particular, no feasible mathematical \(k=2\) candidate with
\(Q_\sigma>0\), physical or spurious, has action below \(A_2(P,Q)\).
\end{theorem}

\begin{proof}
For any feasible broad word, put
\[
 s=h_{D_P}(u),\qquad t=h_{D_Q}(v).
\]
Lemma~\ref{lem:product-two-bounce-block-collapse} gives
\(s,t>0\) and \(s+t=1\).  Therefore
\[
 x=\frac{u}{s}\in\partial D_P^\circ,\qquad
 y=\frac{v}{t}\in\partial D_Q^\circ,
\]
and
\[
 Q_\sigma(\beta)
 =2st\langle x,y\rangle
 \leq2st\,M(P,Q)
 \leq\frac{M(P,Q)}2.
\]

Conversely, compactness supplies maximizers in
\(D_P^\circ\times D_Q^\circ\), and the maximum is positive.  Radially
enlarging either member of a positive maximizing pair shows that both belong
to the respective boundaries.  Thus choose
\(x\in\partial D_P^\circ\) and \(y\in\partial D_Q^\circ\) with
\(\langle x,y\rangle=M(P,Q)>0\), and set
\[
 u=\frac{x}{2},\qquad v=\frac{y}{2}.
\]
Lemma~\ref{lem:product-two-bounce-radial-encoding} encodes
\(u,-u\) as one- or two-facet \(q\)-blocks and \(v,-v\) as one- or
two-facet \(p\)-blocks.  The total coefficient is
\[
 h_{D_P}(u)+h_{D_Q}(v)=\frac12+\frac12=1,
\]
and the four block sums close.  Choosing the cyclic orientation
\(\mathcal Q_+\mathcal P_+\mathcal Q_-\mathcal P_-\) gives
\[
 Q_\sigma=2\langle u,v\rangle=\frac{M(P,Q)}2.
\]
Proposition~\ref{prop:product-two-bounce-distinct-repair} replaces this broad
maximizer by a distinct-facet word of the same value.  Hence
\(Q_{2,\max}=M/2\), and
Lemma~\ref{lem:product-two-bounce-polar-pairing} proves
\eqref{eq:product-two-bounce-complete-qp}.  The project action of a positive
QP candidate is \(1/(2Q_\sigma)\), so the upper bound on \(Q_\sigma\)
also excludes a smaller action from a nonphysical candidate.
\end{proof}
\end{unverified}

\subsection{Physical placement is a separate claim}
\label{subsec:product-two-bounce-physical-placement}

Theorem~\ref{thm:product-strong-two-bounce} constructs a mixed-normal strong
billiard at every minimizing displacement.  After reversing the loop to pass
from Rudolf's Reeb orientation to the repository orientation, splitting a
polygonal normal cone into one or two extremal facet directions yields the
block supports used in the finite QP.  Conversely, reversing back after
summing a physically placed factor-pure generalized characteristic with two
blocks of each type gives the four strong relations
\eqref{eq:product-two-bounce-strong-relations}.

This does not make every feasible or stationary QP word physical.  A feasible
weight vector first gives a closed dual curve.  A translation placing all its
pieces on the prescribed facets is an additional base-point condition, as
recorded in \texttt{formal/reeb-orbit-recovery.tex}.  The global Haim--Kislev
maximizer has such a placement through dual minimality and Clarke
reconstruction, but a \(k=2\) class maximizer need not be the global capacity
maximizer when a three-bounce trajectory is shorter.  Thus the global
reconstruction argument cannot be restricted merely by replacing the full
word family with the \(k=2\) family.

The complete-QP theorem deliberately needs no physicality claim.  Its polar
bound applies to every algebraically feasible word, including spurious ones,
and its broad equality construction followed by the affine repeated-facet
repair proves equality inside the complete finite family.

\subsection{Transition-pruned and exactly aggregated stream}
\label{subsec:product-two-bounce-stream}

The current implementation has further layers:
\begin{enumerate}
  \item
    \texttt{crates/symplectic/src/algorithms/billiard/%
    block\_enumeration.rs} generates distinct one- or two-facet blocks and
    their \(k=2,3\) alternating words.
  \item
    \texttt{crates/symplectic/src/algorithms/billiard/mod.rs} discards words
    failing \texttt{is\_feasible\_cycle}; its transition matrix combines
    facet intersection with the directed sign of \(\omega_0\).
  \item
    \texttt{solve\_billiard\_candidates} sends the surviving words through
    the f64 KKT path.
  \item
    \texttt{class-minima.rs} groups the resulting payloads by bounce count
    and calls
    \path{aggregate_certified_orbits_with_dual_vertices_exact}.
\end{enumerate}
For each candidate that the exact aggregator resolves,
\texttt{orbit\_search.rs} checks exact positivity of all weights, exact
normalization, exact closure in all four coordinates, and positivity of the
exact \(Q\), then records the exact action \(1/(2Q)\).  Its API explicitly
states that the scalar is a global capacity only under candidate-family
coverage and that the candidates are not automatically physical orbits.
Exact aggregation therefore certifies its resolved inputs; it does not prove
that transition pruning or the f64 frontend omitted no relevant word.

\begin{unverified}
\begin{proposition}[Sufficient condition for the stored \(k=2\) value]
\label{prop:product-two-bounce-stream-condition}
For a rational input pair \(P,Q\), suppose the final exactly certified
\(k=2\) candidate set used by \texttt{class-minima.rs} contains a candidate
with
\[
 Q=Q_{2,\max}(P,Q)=\frac1{2A_2(P,Q)}.
\]
Then its stored two-bounce class minimum is exactly \(A_2(P,Q)\).
\end{proposition}

\begin{proof}
The displayed candidate supplies stored action \(A_2(P,Q)\).  Every exactly
certified stream candidate is a feasible mathematical \(k=2\) word, so
Theorem~\ref{thm:product-two-bounce-complete-qp} gives it action at least
\(A_2(P,Q)\), whether or not it is physical.  The minimum over the certified
stream is therefore \(A_2(P,Q)\).
\end{proof}
\end{unverified}

The smallest unresolved universal implementation condition is precisely the
hypothesis of
Proposition~\ref{prop:product-two-bounce-stream-condition}: at least one
complete-family equality word must survive the directed transition filter,
be accepted by the f64 proposal path, and be resolved by the exact
aggregation.  The available source does not prove this.  The transition tests
are necessary local tests, not a sufficient base-point theorem.  The affine
repeated-facet repair preserves closure, normalization, and objective, but no
already established primal base point.  Conversely, the mixed-normal strong
billiard does not by itself select a distinct pure-facet ordering proved to
pass all four directed transitions.  This is the exact remaining gap.

There is no separate risk of an exactly certified nonphysical \(k=2\)
candidate lowering the stored value: the polar upper bound in
Theorem~\ref{thm:product-two-bounce-complete-qp} excludes that algebraically.
A coverage failure can omit all equality candidates and thus produce a larger
value or no \(k=2\) class.

\subsection{Scope of the retained exact comparison}
\label{subsec:product-two-bounce-retained-evidence}

The width-shortcut packet in
\begin{center}
\path{experiments/sys-datascience/methods/product-bounce-width-shortcut/}
\end{center}
is not part of the reviewed proof commit
\texttt{995060e89d4d321b15c786e9391e31d7883de5e9}.  Its current source is
the sibling commit
\texttt{ff636417baa68013033a11a3fc411fb7145515e5}, which has not been
imported or merged here.  On that sibling commit, the packet
computed \(A_2(P,Q)\) independently from factor vertices on 20 deterministic
retained rows, two from each retained \((k,m)\)-bucket.  Its
\texttt{artifacts/summary.json} records 20 exact matches and zero exact
mismatches against the stored two-bounce class minima.

This is finite validation of those 20 rows only.  It does not prove the
universal stream-coverage condition, physicality of the reported minimizing
words, completeness on the other 10,220 retained rows, or absence of f64
false negatives on other inputs.  It also says nothing about the
three-bounce class or which class realizes the global EHZ capacity.  The
geometric and complete-QP theorems above are exact arguments independent of
these comparisons.

\subsection{Proved and unproved boundary}
\label{subsec:product-two-bounce-conclusion}

At agent-review status, the packet proves
\[
 \boxed{
 \min_{\text{strong two-bounce billiards}}\ell_Q
 =\min_{d\in\partial(P-P)}h_{Q-Q}(d)
 =\max\{r:r(Q-Q)^\circ\subseteq P-P\}
 }
\]
and
\[
 \boxed{
 Q_{2,\max}^{\mathrm{complete}}
 =\frac12
   \max_{\substack{x\in(P-P)^\circ\\y\in(Q-Q)^\circ}}
   \langle x,y\rangle
 =\frac1{2A_2(P,Q)}.
 }
\]
The first identity includes strong realization of every minimizing
displacement.  The second includes the distinct-facet affine repair and rules
out smaller spurious \(k=2\) QP actions.

Unproved universally is equality with the actual transition-pruned,
f64-seeded stream.  It follows under the single explicit coverage condition
in Proposition~\ref{prop:product-two-bounce-stream-condition}.  The 20/20
retained exact comparison validates only those rows.  Agent review and a
successful build do not constitute Jörn's mathematical acceptance.
